\chapter{全文总结}
\label{chap:conclusion}

本文以“水下发射机—海面漂浮液晶 RIS—空中接收机”构成的水空直接无线光链路为研究对象，按照海面平台运动、RIS 下表面入射光场、RIS 相位设计与接收性能的顺序建立了系统模型，并以数值仿真评估了液晶 RIS 对链路性能的增强作用。主要工作与结论如下。

（1）建立了海面漂浮液晶 RIS 的器件方案与六自由度运动模型。以 $40\times40$ 个 14.4\,mm 通光区的液晶单元拼接构成 $1\,\mathrm m\times1\,\mathrm m$ 玻璃夹层面板，1\,$\mu$m 的微电极周期给出 $15.426^\circ$ 的单轴偏转角上限和 $0.021^\circ$ 的量化步长。RIS 下层玻璃贴水安装时，水—玻璃界面不发生全反射，液晶层在光束进入空气之前减小其切向波矢，单轴条件下全反射的起始入射角由 $48.75^\circ$ 推迟到 $72.15^\circ$。完成了 29.871\,kg 整机的浮框配平，静水中玻璃浸水 3.432\,mm、顶面高出水面 6.568\,mm。以 Neumann 方向谱和梯形半权重离散生成随机海面，默认网格下离散方差误差不超过 0.133\%；建立了四元数描述的六自由度刚体方程，耦合薄板辐射/绕射载荷、侧壁压力与解析偏航力矩以及被动辐射记忆模型，得到了风速 3--12\,m/s 下 RIS 由厘米级随机晃动到米级长波轨道运动的姿态与轨迹。

（2）建立并验证了 RIS 下表面入射光场的光束扩展函数模型。把入射光强分解为弹道分量与散射分量：弹道分量对有限束腰源在任意 RIS 姿态下给出精确的面积 Jacobian 表达式和闭式点源极限；散射分量以完整 HG 相函数下辐射输运的碰撞次数展开为定义，以强制首次散射估计器无偏求值，并把整体二维数值标准误差控制在 0.1\% 以内。在由运动模型冻结的 18 个快照上，模型与独立 MC 的差异与 MC 自身的统计噪声一致，总功率差不超过 0.317\%。小角近似 S-BSF 在入射角不超过约 $65^\circ$ 时与完整模型的偏差小于 0.4\%，在随风浪变化的掠入射场景中偏差增大到 2\%--3\%。在此基础上给出了液晶层上的入射复光场与单元功率。

（3）提出了受器件约束的有限宽度光斑整形相位设计方法，建立了从 RIS 到接收端的完整性能评估模型。以入射能量的白化径向分位数为单元分配目标落点，比较平顶、高斯和径向 Weibull 型鲁棒三类目标分布，经分层 Fermat 路径反解、逐轴限幅量化和实际落点上的常相位校准得到可加载相位，并以半能量半径约束避免聚焦为单点。采用二阶光程展开下的逐单元衍射积分、CZT、Parseval 功率守恒搬运与像素内相干平均计算接收面光强，联合计及瑞利位置抖动、Rice 角度误差、倾斜有限孔径、有限视场和 EGG--Gamma-Gamma 联合衰落，以平均功率、ABER 和 OP 三类目标选择相位。

（4）系统评估了 RIS 相位方案对链路性能的影响。液晶 RIS 把 $40^\circ$ 入射时光斑质心的偏离角由 $58.8^\circ$ 降到 $29.3^\circ$--$32.7^\circ$，消除斜入射像散，10\,m 处峰值光强达到无 RIS 的 10.3--12.6 倍，代价是 33.2\% 填充率带来的总功率损失。共同固定接收机下，相位调制相对未调相 RIS 带来约 40\,dB 的功率增益。对准自身光斑质心的直接传播在入射角低于约 $42^\circ$--$44^\circ$ 时性能更优，高于该角度时 RIS 鲁棒方案更优，超过临界角后 RIS 是唯一可用的上行通道。鲁棒光斑在大位置抖动下优于最大平均功率方案和平顶方案：$\sigma_p=120\,\mathrm{mm}$、45\,dBm 时，最大平均功率方案的平均接收功率为鲁棒方案的 1.68 倍，平均误码率却是鲁棒方案的 9.9 倍。最优光斑半径随位置抖动标准差近似线性增大，约为其 2.4--2.9 倍。位置抖动的影响取决于光斑形状，可以由相位设计缓解；角度误差主要通过视场失收形成与光斑形状无关的误码平台，需要依靠接收端的视场与指向设计解决。

% 【预留】后续工作展望一节待动态海况研究推进后补充。
% \section{后续工作展望}
